\scanpage{71}
\mspage{15}

Par contre, ce n'est nullement le cas en prenant $\unsure{\mathcal F|\mathscr F}=\underline{\mathrm{Fib}}\,S,\underline{\mathrm{Liss}}\,S$ ou $\underline{\mathrm{Cofib}}\,S,\underline{\mathrm{Prop}}\,S$. Ayant choisi dans ces catégories comme localisateurs \unsure{la notion|les notions} de $W$-équivalences ``fibre par fibre'', j'ai donc l'air de dire que je trouve pour ces quatre catégories les \uline{mêmes} \unsure{localisées|localisés} (à équivalence près) que si je prenais le localisateur beaucoup plus grossier hérité par $\mathrm{Cat}_{/S}$. Cela semble difficile à croire, et il faudra que je fasse très attention en démontrant les énoncés-clef.

Pour voir que les flèches en traits pleins du grand diagramme p.~6 sont toutes compatibles avec les localis(at)eurs, donc aussi les flèches composées au moyen de celles-ci, j'ai à utiliser déjà ceci (pour \inserted{le cas} \deleted{\unsure{de}}des flèches $\underline{\mathrm{Liss}}\,S\unsure{\text{---}|\to}\mathrm{Cat}_{/S},\ \underline{\mathrm{Prop}}_{/S}\to\mathrm{Cat}_{/S}$)\,:

\textcircled{1} Si $X,Y$ lisses \inserted{ou propres} sur $S$, et $f:X\unsure{\text{---}|\to}Y$ un $S$-\deleted{\unsure{fonc}}\unsure{foncteur}, si $f_s:X_s\to Y_s$ est dans $W$ pour \unsure{tout} $s\in S$, alors $f\in W$.
% slanted
\margin{OK si on suppose L 3 gauche et L 4}

Cela renforce considérablement la condition L 5 sur $W$, où on supposait \deleted{\unsure{que}} que $f$ est un foncteur \uline{cartésien} entre catégories \uline{fibrées} sur $S$.

(À \unsure{dégager|déduire}, pour les différentes propriétés dont on a besoin\unsure{$\ne$|,} pour $W$, à quelles propriétés \unsure{standart|standard} on peut les ramener. Tout sera OK si $W=W_\infty$.)

% archive sheet number: 70
